\section{Como a Solução Funciona}

O validador lê a placa da esquerda para a direita, um caractere
por vez, e mantém apenas uma informação: em que ``etapa'' da placa
está. Por exemplo, depois de ler três letras, a etapa seguinte
espera um algarismo (ou um hífen, no padrão anterior). Se aparecer
um caractere que não cabe naquela etapa, a placa é descartada na
hora. Se a leitura terminar exatamente na última etapa prevista, a
placa é aceita.

Cada ``etapa'' é um \emph{estado} do autômato, e a regra de passagem
de uma etapa para outra é a \emph{função de transição}. Como o
número de etapas é pequeno e fixo, o validador é rápido e seu
comportamento é totalmente previsível: a mesma entrada sempre
produz a mesma resposta.